\subsection{Limitations of the Thermal Guard}
\label{sec:guard-limit}

Samsung uses a static thermal guard to prevent Capture Timeout violations in deployed single-frame draft stages.
The guard uses Samsung's overheat level, an ordinal thermal indicator on a 0--6 scale, where larger values denote stronger throttling.
During pre-release validation, throttling at level~4 repeatedly delayed a single-frame draft stage and caused Capture Timeout violations despite optimization.
The guard therefore skips optional single-frame stages at level~4 or higher.

However, the overheat level does not indicate a capture's waiting time or remaining draft work.
The guard may therefore skip optional stages even when enough time remains to complete the draft sequence.
For example, ambient heat, charging, or prior heavy device workloads can readily raise the overheat level to 4 before a camera session starts.
The guard then skips optional stages on the session's first capture, even when its remaining timeout budget is sufficient.

\smallskip\noindent\textbf{Motivating experiment.}
To examine the timeout risk of enabling lightweight multi-frame composition in the draft sequence, we compared two Portrait configurations on the Galaxy S26 Ultra.
We bypassed the thermal guard so that configured optional stages could execute regardless of the overheat level.
The optional stages are the multi-frame composition stage (\(M\)) and the deployed single-frame stages (collectively \(S\)).
Both \(M\) alone and \(M+S\) retain the same mandatory encoding and saving stages.
We tested 12MP and 24MP capture under normal memory conditions and under memory pressure.
For each condition in Table~\ref{tab:timeout_index}, we ran ten trials of up to 30 consecutive captures.
We chose 30 captures to extend testing beyond the earlier level-4 validation failure near the 20th consecutive capture.
The table reports the earliest timeout capture index across the ten trials, grouped by starting overheat level.

% Notes: docs/exhibits.md#tab_timeout_index
\providecolor{tmoS1}{RGB}{252,227,220} %
\providecolor{tmoS2}{RGB}{239,160,140} %
\providecolor{tmoS3}{RGB}{201,80,63}  %
\providecolor{tmoS4}{RGB}{143,36,32}  %

\newcolumntype{H}{>{\centering\arraybackslash}m{1.20cm}}
\newcolumntype{P}{>{\centering\arraybackslash}m{0.67cm}}

\providecommand{\tmoFirst}[2]{\cellcolor{#1}\makebox[\linewidth][r]{#2\hspace{1pt}}}
\providecommand{\tmoFirstW}[2]{\cellcolor{#1}\makebox[\linewidth][r]{\textcolor{white}{\textbf{#2}}\hspace{1pt}}}
\providecommand{\tmoNone}{\makebox[\linewidth][r]{--\hspace{1pt}}}

\begin{table}[t]
 \centering
 \caption{Capture index of the earliest timeout.}
 \label{tab:timeout_index}
 {\scriptsize\setlength{\tabcolsep}{1.3pt}%
 \begin{tabular}{@{}HPP|PP||PP|PP@{}}
 \toprule
 \multirow[c]{3}{=}[-10pt]{\makecell[c]{\textbf{Starting}\\\textbf{overheat}\\\textbf{level}}} &
 \multicolumn{4}{c||}{\raisebox{-3pt}{\textbf{Normal capture}}} &
 \multicolumn{4}{c@{}}{\raisebox{-3pt}{\textbf{Memory-pressure capture}}} \\[6pt] \cmidrule(lr){2-5}\cmidrule(l){6-9}
 &
 \multicolumn{2}{c|}{\raisebox{-3pt}{\textbf{12MP}}} &
 \multicolumn{2}{c||}{\raisebox{-3pt}{\textbf{24MP}}} &
 \multicolumn{2}{c|}{\raisebox{-3pt}{\textbf{12MP}}} &
 \multicolumn{2}{c@{}}{\raisebox{-3pt}{\textbf{24MP}}} \\[6pt] \cmidrule(lr){2-3}\cmidrule(lr){4-5}\cmidrule(lr){6-7}\cmidrule(l){8-9}
 & \raisebox{-2pt}{\textbf{M{+}S}} & \raisebox{-2pt}{\textbf{M}} &
 \raisebox{-2pt}{\textbf{M{+}S}} & \raisebox{-2pt}{\textbf{M}} &
 \raisebox{-2pt}{\textbf{M{+}S}} & \raisebox{-2pt}{\textbf{M}} &
 \raisebox{-2pt}{\textbf{M{+}S}} & \raisebox{-2pt}{\textbf{M}} \\[4pt]
 \midrule
 Lv0 &
 \tmoNone & \tmoNone & \tmoFirst{tmoS1}{29} & \tmoNone &
 \tmoNone & \tmoNone & \tmoFirst{tmoS1}{26} & \tmoNone \\
 Lv1 &
 \tmoNone & \tmoNone & \tmoFirst{tmoS1}{29} & \tmoNone &
 \tmoFirst{tmoS1}{24} & \tmoNone & \tmoFirst{tmoS2}{19} & \tmoNone \\
 Lv2 &
 \tmoFirst{tmoS1}{24} & \tmoNone & \tmoFirst{tmoS2}{18} & \tmoNone &
 \tmoFirst{tmoS2}{15} & \tmoFirst{tmoS1}{29} &
 \tmoFirst{tmoS2}{14} & \tmoFirst{tmoS1}{25} \\
 Lv3 &
 \tmoFirst{tmoS2}{13} & \tmoFirst{tmoS1}{21} &
 \tmoFirst{tmoS2}{12} & \tmoFirst{tmoS1}{20} &
 \tmoFirstW{tmoS3}{8} & \tmoFirst{tmoS2}{19} &
 \tmoFirstW{tmoS4}{5} & \tmoFirst{tmoS2}{13} \\
 \midrule
 Lv4 &
 \tmoFirstW{tmoS3}{8} & \tmoFirst{tmoS2}{16} &
 \tmoFirstW{tmoS3}{7} & \tmoFirst{tmoS2}{13} &
 \tmoFirstW{tmoS3}{6} & \tmoFirst{tmoS2}{13} &
 \tmoFirstW{tmoS4}{3} & \tmoFirstW{tmoS3}{7} \\
 Lv5 &
 \tmoFirstW{tmoS3}{8} & \tmoFirst{tmoS2}{12} &
 \tmoFirstW{tmoS3}{8} & \tmoFirst{tmoS2}{12} &
 \tmoFirstW{tmoS4}{4} & \tmoFirstW{tmoS3}{9} &
 \tmoFirstW{tmoS4}{4} & \tmoFirstW{tmoS3}{9} \\
 Lv6 &
 \tmoFirstW{tmoS3}{7} & \tmoFirst{tmoS2}{11} &
 \tmoFirstW{tmoS3}{7} & \tmoFirst{tmoS2}{11} &
 \tmoFirstW{tmoS4}{3} & \tmoFirstW{tmoS3}{6} &
 \tmoFirstW{tmoS4}{3} & \tmoFirstW{tmoS3}{6} \\
 \bottomrule
 \end{tabular}%
 }
 \par\vspace{3pt}
 {\scriptsize\begin{minipage}{0.933\columnwidth}
 \raggedright\setlength{\fboxsep}{1.8pt}%
 Timeout capture index:~%
 \colorbox{tmoS1}{\strut\,${\geq}\,20$\,}\colorbox{tmoS2}{\strut\,10--19\,}\colorbox{tmoS3}{\strut\,\textcolor{white}{6--9}\,}\colorbox{tmoS4}{\strut\,\textcolor{white}{${\leq}\,5$}\,}\par
 --: no timeout observed within 30 captures.\par
 \end{minipage}}
\end{table}


The overheat level alone was insufficient to explain timeout onset.
With \(M+S\), timeouts occurred even in runs starting below the guard's cutoff.
Conversely, every \(M+S\) trial starting at level~4 or higher completed its first capture without a timeout.
At the same starting level on the same device, the earliest timeout varied with resolution, memory pressure, and draft configuration.
In particular, compared with \(M+S\), using \(M\) alone delayed the earliest timeout or resulted in no observed timeout within 30 captures.

\smallskip\noindent\textbf{Need for runtime control.}
The experiment shows that reducing optional draft work can delay timeout onset.
However, skipping optional stages can widen the draft--final visual gap and cannot recover timeout budget already spent waiting.
Delaying the next capture opportunity can limit draft queue growth and preserve budget for optional processing, at the cost of a longer shot-to-shot interval.
These trade-offs motivate coordinated runtime control of optional-stage execution and capture arrivals.
To determine how much timeout budget to reserve, both decisions require estimating the time needed for the remaining draft work.
Reserving too little risks a Capture Timeout, whereas reserving too much needlessly sacrifices draft image quality or capture responsiveness.
